% ----------------------------------------------------------
\chapter{Considerações Finais}
% ----------------------------------------------------------

O presente trabalho cumpriu com seus objetivos, demonstrando a viabilidade de desenvolver um protótipo de modelador de diagramas ER colaborativo acessível via navegador. A plataforma resultante integra, em uma única solução gratuita, a modelagem interativa de esquemas de banco de dados, a colaboração em tempo real entre múltiplos usuários e operações de importação e exportação DDL, conjunto de capacidades que, conforme evidenciado na comparação com ferramentas do mercado, não encontra equivalente gratuito disponível.

Os cinco objetivos específicos definidos na introdução foram atendidos. O primeiro objetivo, referente à modelagem e persistência dos elementos do diagrama, foi atendido pelo editor interativo desenvolvido, que captura entidades, atributos, relacionamentos, tipos de dados e posicionamentos, preservando o estado completo do diagrama. O segundo objetivo, de investigar a principal biblioteca de diagramação adotada, foi cumprido ao longo do desenvolvimento: a biblioteca demonstrou-se adequada para a renderização interativa de diagramas ER e para a aplicação de atualizações provenientes de outros colaboradores em tempo real. O terceiro objetivo, de permitir a colaboração simultânea de múltiplos usuários, foi alcançado por meio de comunicação bidirecional em tempo real, com mecanismo de controle de acesso por entidade para prevenção de conflitos. O quarto objetivo, de implementar autenticação e autorização, foi satisfeito pela combinação de mecanismos de segurança no servidor e no cliente. O quinto objetivo, de adotar práticas ágeis, foi contemplado pela aplicação adaptada de metodologia ágil com acompanhamento visual do progresso durante todo o ciclo de desenvolvimento.

O ganho proporcionado por este trabalho se manifesta em duas dimensões. Do ponto de vista técnico, o protótipo serve como referência prática para a construção de ferramentas colaborativas baseadas em comunicação em tempo real com editores visuais interativos em ambiente \textit{web}. Do ponto de vista da área, o trabalho contribui para preencher a lacuna identificada na justificativa: a carência de uma ferramenta colaborativa, baseada em \textit{web} e sem barreiras de custo para a modelagem de diagramas ER, atendendo às demandas de equipes de desenvolvimento distribuídas.
